\subsection{Evaluation design}
\label{sec:mujoco_3d}

The primary comparison evaluated unchanged BAMOR-v25 and fully seeded CAPQL
at 75k interactions on HalfCheetah and Hopper, using seeds 1601--1620 and only
the final actor. An earlier five-seed pilot motivated the task selection and
utility hypothesis; confirmation seeds and analysis were fixed before training.
The cohorts remain separate (Appendix~\ref{app:experiment_details}).

The custom piecewise-stationary tasks have two objectives/four modes on
HalfCheetah, three/four on Hopper and three/three on Ant. Modes alter physics
and reward scales or smooth caps; exponential dwell lengths have scale 25k
and minimum 1,000. Evaluation locks each mode, uses deterministic actions at
21 full-simplex weights with three paired reset seeds, and ends on termination
or at 1,000 steps. Training schedules and evaluation inputs are paired;
mode-wise fronts are not pooled (Appendix~\ref{app:task_details}).

Direct utility $U$ evaluates the requested preference; final HV evaluates the
final actor's return set; archive HV evaluates retained actors. W/L/T denotes
paired seed wins/losses/ties, including every pair. Relative differences compare
cohort means; utilities and raw HV are not pooled across tasks. BAMOR uses width
64 and CAPQL width 256, with equal interaction budgets but different update
computation (Table~\ref{tab:v23_budget}).

\subsection{Twenty-seed final-policy utility}
\label{sec:utility_confirmation}

BAMOR achieved higher observed mean direct utility on both tasks
(Table~\ref{tab:utility_confirmation}). HalfCheetah means were 296.32 versus
191.17 for CAPQL, a difference of 105.15 (+55.00\%); Hopper means were 436.59
versus 411.93, a difference of 24.67 (+5.99\%). Each task had 13 positive
and seven negative seed-paired differences. Every run used the original
method settings, with no intermediate evaluation or checkpoint selection.

% Generated by make_evidence_tables.py from completed readouts.
\begin{table}[t]
\centering\small
\caption{Final-policy direct utility at 75k training interactions, 20 independent seed pairs per environment (1601--1620). Means are BAMOR / CAPQL; $\Delta U$ and its interval are in utility units. The paired 97.5\% bootstrap intervals implement the prespecified two-environment criterion. W/L/T includes every seed pair.}
\label{tab:utility_confirmation}
\begin{tabular}{@{}lrrrrr@{}}
\toprule
Environment & $U$: BAMOR / CAPQL & $\Delta U$ & Relative & 97.5\% CI & W/L/T \\
\midrule
HalfCheetah & 296.32 / 191.17 & +105.15 & +55.00\% & $[-33.80,\ 235.83]$ & 13/7/0 \\
Hopper & 436.59 / 411.93 & +24.67 & +5.99\% & $[-58.64,\ 105.58]$ & 13/7/0 \\
\bottomrule
\end{tabular}
\end{table}


\paragraph{Prespecified uncertainty analysis.}
Intervals used 20,000 bootstrap resamples of complete training-seed pairs,
retaining their mode and preference outcomes together. The frozen criterion
used a 97.5\% percentile interval per task for nominal 95\% Bonferroni family
coverage; bootstrap coverage is approximate. Both intervals included zero, as
did ordinary 95\% intervals (Appendix~\ref{app:experiment_details}). The
observed means therefore did not meet the criterion for an expected-utility
advantage. No pairs were excluded and no pilot seeds were added.

\paragraph{Coverage and computational cost.}
Final HV and runtime gave different orderings from mean utility
(Table~\ref{tab:utility_secondary}). BAMOR's mean final HV was 4.84\% lower
on HalfCheetah and 5.77\% lower on Hopper, with 10 and eight wins out of 20,
respectively. The descriptive paired HV intervals included zero. Mean
training time was 1.67 and 1.73 times CAPQL's. These are measured times on a
shared cluster, not dedicated-machine benchmarks; both methods used two CPU
threads per run.

% Generated by make_evidence_tables.py from completed readouts.
\begin{table}[t]
\centering\small
\caption{Secondary outcomes on the same 20-seed cohorts as Table~\ref{tab:utility_confirmation}. HV is the relative difference of final-policy means. Times are mean seconds per run, BAMOR / CAPQL, on the shared CPU cluster with two threads per run.}
\label{tab:utility_secondary}
\begin{tabular}{@{}lrrrr@{}}
\toprule
Environment & Final HV & W/L/T & Training (s) & Evaluation (s) \\
\midrule
HalfCheetah & -4.84\% & 10/10/0 & 1635.8 / 981.5 & 152.5 / 81.8 \\
Hopper & -5.77\% & 8/12/0 & 1759.7 / 1016.8 & 122.6 / 69.1 \\
\bottomrule
\end{tabular}
\end{table}


The 80 runs consumed 6,000,000 training and 16,512,980 evaluation interactions.
Evaluation did not enter replay or selector feedback. Initialization probes,
prelaunch tests and per-cohort costs are accounted for in
Appendix~\ref{app:experiment_details}.

\subsection{Archive coverage and feedback attribution}

At 50k interactions, both methods retained 19 actors. This retrospective
comparison used 15 HalfCheetah and five Hopper/Ant seeds, adding CAPQL after
the BAMOR cohort was inspected. Archive coverage favored BAMOR most strongly
on Hopper, with five paired wins (Table~\ref{tab:matched_archive}). Ant final
HV fell on all five seeds; HalfCheetah's 16.21\% archive-mean gain accompanied
only 0.74\% in final HV. The early CAPQL cohort's incomplete RNG seeding is
documented in Section~\ref{sec:limits}.

% Generated by make_evidence_tables.py from completed readouts.
\begin{table}[t]
\centering\small
\caption{Retrospective 50k comparison: v25 minus CAPQL, with 19 actors in each archive. Percentages compare cohort means; W/L/T denotes paired seed wins/losses/ties. The early CAPQL runs omitted two RNG seeds (Section~\ref{sec:limits}).}
\label{tab:matched_archive}
\begin{tabular}{lrrrrr}
\toprule
Environment & $n$ & Archive HV & W/L/T & Final HV & W/L/T \\
\midrule
HalfCheetah & 15 & +16.21\% & 10/5/0 & +0.74\% & 8/7/0 \\
Hopper & 5 & +140.39\% & 5/0/0 & +58.67\% & 3/2/0 \\
Ant & 5 & +1.17\% & 3/2/0 & -11.07\% & 0/5/0 \\
\bottomrule
\end{tabular}
\end{table}


The score-blind control retained networks, arms, family floors and archive
schedule while replacing routing scores with zero
(Table~\ref{tab:selector_effect}). Adaptive selection did not improve mean
archive or final HV on HalfCheetah or Ant. Hopper's positive means had two
wins, one loss and two ties; seven of fifteen HalfCheetah pairs tied exactly.
These observations separate an external archive advantage from evidence of a
consistent feedback benefit.

% Generated by make_evidence_tables.py from completed readouts.
\begin{table}[t]
\centering\small
\caption{Feedback attribution at 50k: unchanged v25 adaptive minus its score-blind control on the same cohorts as Table~\ref{tab:matched_archive}.}
\label{tab:selector_effect}
\begin{tabular}{lrrrrr}
\toprule
Environment & $n$ & Archive HV & W/L/T & Final HV & W/L/T \\
\midrule
HalfCheetah & 15 & -0.15\% & 3/5/7 & -0.85\% & 4/4/7 \\
Hopper & 5 & +6.54\% & 2/1/2 & +59.56\% & 2/1/2 \\
Ant & 5 & -1.00\% & 1/3/1 & -3.33\% & 0/4/1 \\
\bottomrule
\end{tabular}
\end{table}


\subsection{Ant mixed-preference diagnosis}
\label{sec:ant_diagnosis}

The earlier fully seeded 75k pilot lost all five Ant pairs in utility and final
HV: mean utility was 988.8 versus CAPQL's 1364.9, and mean final HV was 6.94\%
lower (Table~\ref{tab:final_policy}). Failure was concentrated in mixed weights
containing the energy objective. On that inspected cohort, re-scalarizing evaluated
returns showed that v25's pure-energy policy beat its corresponding
mixed-preference policy at 66 of 70 mode-0 seed--weight cells. This comparison
uses the same saved actor and requires no extra environment interactions.

Candidate-search ablations and an independent structure-by-selector pilot did
not close the CAPQL gap. The latter's primary residual effect had only two of
five wins, and its residual adaptive arm still trailed CAPQL by 247.9 mixed
utility and 2.06\% final HV (Appendix~\ref{app:ant_diagnostics}). The residual
variant was not adopted; cross-environment BAMOR remains unchanged v25.
